\section{Conclusion}
\label{sec:conclusion}

We composed two existing universal-attack methods --- the Universal Adversarial Attack of \citet{rahmatullaev2025universal} and the \textsc{AnyAttack} encoder--decoder of \citet{zhang2025anyattack} --- with a programmatic dual-axis evaluation that scores Output Affected (drift) and Target Injected (payload delivery) independently. Across $6{,}615$ (clean, adversarial) response pairs over four open \vlm{}s, seven attack prompts, and seven test images at $\Linf = 16/255$ and PSNR $\approx 25.2$~dB, broad disruption ($\approx 66\%$) coexists with rare literal injection ($0.227\%$) --- a $290\times$ divergence on the same data. The few injections that do land cluster on screenshot-like images whose semantic class matches the payload's category; a manual transfer test against GPT-4o fails; BLIP-2 is fully immune across every pair we evaluate, even when used as a Stage-1 surrogate.

Two implications stand out. First, single-number ``ASR'' metrics in the universal-attack-on-VLM literature conflate disruption with payload delivery, and the two diverge by orders of magnitude. Future evaluations should report both axes. Second, the BLIP-2 result suggests that information-bottleneck architectures (Q-Former-style) may already provide substantial adversarial robustness against \eps{}-bounded attacks at no additional training cost --- a defense direction that does not require expensive adversarial fine-tuning of the vision encoder.

\paragraph{Future work.}
Three natural extensions are out of scope here. (i) A direct $x_u \to$ BLIP-2 ablation (skipping Stage 2 fusion) would distinguish whether the immunity sits in the AnyAttack decoder or in BLIP-2's architecture itself. (ii) A systematic transferability study across multiple frontier closed models (ChatGPT, Gemini, Claude) on multiple cases would replace this paper's single GPT-4o anecdote with a measured rate. (iii) Extending the attack family beyond gradient-based pixel perturbation (typographic injection, steganographic embedding, scene spoofing) would test whether the disruption-versus-injection gap holds across the wider visual-injection design space.

\paragraph{Reproducibility.}
All experiments use deterministic seeds where applicable; both stages of the attack use public pretrained weights without retraining (\texttt{coco\_bi.pt} for AnyAttack); the dual-axis judge contains no LLM in the loop. Stage~1 takes $7$--$19$ minutes per universal image on a single H200 80~GB; Stage~3 generation is the dominant wall-clock cost ($\sim 30$ minutes per (image, experiment) row); Stage~3 evaluation runs in ${\sim}5$ minutes on a laptop CPU. Hyperparameters: PGD steps $= 2{,}000$, learning rate $10^{-2}$ (Adam), $\eps = 16/255$, ensemble sizes $N \in \{2, 3, 4\}$. Full training and evaluation parameters are listed in the released code at \url{github.com/jeffliulab/vis-inject}.

\paragraph{Released artifacts.}
\begin{itemize}
  \item Dataset (CC-BY-4.0): \url{huggingface.co/datasets/jeffliulab/visinject} --- 21 universal adversarial images, 147 adversarial photos, $6{,}615$ response pairs with dual-axis judge scores, and 12 curated case studies.
  \item Code (MIT license): \url{github.com/jeffliulab/vis-inject} --- full pipeline, all VLM wrappers, and the dual-axis judge.
  \item Interactive demo: \url{huggingface.co/spaces/jeffliulab/visinject}.
\end{itemize}
